\section{Contexto y justificación}

\begin{nota}{Centro ficticio.}
El IES Ejemplo, su grupo-clase y sus documentos son una descripción inventada, verosímil y genérica, creada para contextualizar esta propuesta en un instituto público de la Comunidad de Madrid.
\end{nota}

\subsection{El centro}
El IES Ejemplo es un instituto público de un municipio del área metropolitana de Madrid (latitud aproximada de 40,4°~N), con unos 800 estudiantes de ESO y Bachillerato de las modalidades de Ciencias y Tecnología y de Humanidades y Ciencias Sociales. Tiene un edificio principal de tres plantas, un gimnasio exento y un patio con pistas deportivas. Al sur del patio, tras la valla del centro, se levanta un bloque de viviendas a cuyo pie no se puede llegar. Hay dos carros de portátiles y conexión inalámbrica en todas las aulas. En la cubierta del gimnasio funciona una instalación fotovoltaica de autoconsumo, y el equipo directivo estudia ampliarla.

\subsection{El grupo-clase}
El grupo de 1.º de Bachillerato A cuenta con 28 estudiantes de 16 y 17 años. Es un grupo participativo y con niveles heterogéneos en Matemáticas: un sector numeroso domina los procedimientos pero tiene dificultad para justificarlos y otro necesita apoyo para pasar del triángulo rectángulo a los problemas generales. Hay un alumno con TDAH, una alumna con dislexia, dos estudiantes de alto rendimiento y una alumna incorporada este curso. Todo ello condiciona las medidas de la sección~\ref{sec:diversidad}.

\subsection{Organización y tutoría}
Matemáticas I es una materia específica de modalidad con cuatro horas lectivas semanales \parencite[anexo~I, p.~32]{decreto64}, en periodos de 55 minutos. El departamento de Matemáticas coordina la propuesta y la comisión de coordinación pedagógica la articula con el resto de materias. El profesor tutor realiza el seguimiento del proceso educativo del grupo \parencite[art.~5.1, p.~17]{decreto64}, y el profesor de Matemáticas le informa del funcionamiento de los equipos. Como el horario del anexo~I no incluye una hora de tutoría en este curso \parencite[p.~32]{decreto64}, la formación en los roles se realiza dentro de las sesiones de Matemáticas\soloProducto{ (sección~\ref{sec:cooperativo})}. El departamento de orientación puede aportar información sobre salidas profesionales vinculadas a la energía, en línea con la finalidad orientadora del plan de acción tutorial \parencite[art.~5.3, p.~17]{decreto64}.

\subsection{Documentos del centro}
\begin{itemize}
  \item \textbf{Proyecto educativo:} recoge valores, fines y prioridades del centro e incorpora la concreción del currículo \parencite[art.~121.1, p.~122921]{lomloe}. En el IES Ejemplo una de sus líneas prioritarias es la sostenibilidad del centro.
  \item \textbf{Programación general anual:} incluye este curso el objetivo de estudiar la ampliación de la instalación solar. El Consejo Escolar la aprueba y evalúa \parencite[art.~127.b, p.~122923]{lomloe}.
  \item \textbf{Plan de acción tutorial:} forma parte de la programación general anual \parencite[art.~5.3, p.~17]{decreto64}.
  \item \textbf{Programación didáctica de Matemáticas I:} sitúa esta propuesta en el segundo trimestre, cuando se abordan geometría y trigonometría.
\end{itemize}

\subsection{Marco legislativo}
La propuesta se rige por la Ley Orgánica 3/2020 \parencite{lomloe}, el Real Decreto 243/2022, que fija las enseñanzas mínimas y los descriptores del perfil de salida \parencite{rd243}, y el Decreto 64/2022 de la Comunidad de Madrid \parencite{decreto64}, cuyo anexo~II recoge las competencias específicas, los criterios y los contenidos de Matemáticas~I. El decreto pide planificar actividades que pongan en contexto la enseñanza y el aprendizaje \parencite[art.~16.3, p.~21]{decreto64} y permite a los centros adoptar innovaciones pedagógicas en el marco de su autonomía \parencite[art.~19.3.a, pp.~21--22]{decreto64}.

\subsection{Justificación de la propuesta}
\textbf{Por qué esta propuesta en este curso.} El currículo destaca el carácter instrumental de las matemáticas en otras áreas científicas y tecnológicas y sitúa la resolución de problemas y la investigación como aspectos fundamentales de la materia \parencite[p.~320]{decreto64}. Vectores, trigonometría y rectas suelen enseñarse por separado y este alumnado tiende a verlos sin relación. El problema de la sombra los necesita a todos: el rayo de sol es una recta con un vector director, su ángulo procede de la trigonometría y las alturas inaccesibles se obtienen resolviendo triángulos.

\textbf{Por qué en este grupo.} El grupo domina los procedimientos, pero le cuesta justificarlos. Defender una decisión real ante quien decide obliga a argumentar, y los equipos heterogéneos aprovechan la diversidad de niveles.

\subsection{Relación con otras materias}
\begin{itemize}
  \item \textbf{Tecnología e Ingeniería I:} energías renovables, eficiencia energética y sostenibilidad \parencite[p.~385]{decreto64}. Aporta los criterios técnicos de la instalación.
  \item \textbf{Física y Química:} carácter vectorial de las magnitudes y composición vectorial \parencite[p.~166]{decreto64}.
  \item \textbf{Orientación:} salidas profesionales vinculadas a la energía.
\end{itemize}

\subsection{Objetivos de Desarrollo Sostenible}
\begin{itemize}
  \item \textbf{ODS~7} \parencite{onu7}: la meta~7.2 propone aumentar la proporción de energía renovable. El alumnado estudia cómo aprovechar mejor una instalación solar real.
  \item \textbf{ODS~11} \parencite{onu11}: la meta~11.6 busca reducir el impacto ambiental de las ciudades. El instituto actúa como pequeña comunidad que planifica el uso de su propio espacio.
\end{itemize}
